\chapter*{The Vision of the Risen Christ}
\textbf{"After that, he was seen of above five hundred brethren at
once."}--\textit{Paul}.


Once for himself was Quintus to see the Lord, before his departure
heavenward.  When midnight hours afterward came to him in Italy,
the memory of that vision was golden.  When, among the temples of
the gods in pagan Rome, men challenged his belief, his sufficient
answer was: "With mine own eyes I have seen the risen Teacher who
has revealed immortality to men."  So did the first disciples of
the faith who bore its weightiest burdens, enjoy its highest
privilege.

It was the disciple John who told Quintus of the opportunity to see
the risen Lord.  In an hour of fellowship at Jerusalem--when the
knight had confessed his new allegiance--John spoke of the Master's
wish.  The disciples who were in the city and its environs were to
gather in Galilee with those from that upper district.  Once more
would their Lord show himself to all who believed on him, and would
speak with them.  Nor did Quintus ever cease to rejoice that he was
reckoned worthy to look that day on the Conqueror of death.

With light feet the Jerusalem company, some six score in number,
made the journey north to Galilee.  One subject only was on their
lips, as they followed the road through Samaria to Kurn Hattin,
near the Sea of Tiberias.  Here the Lord at the opening of his
mission had spoken his nine blessings to needy mortals; most
fitting it now was that on this memorable hillside he should utter
his farewell to those who had come to believe on him.  Thus would
the circle of his teachings end where it had begun.  Bright was the
picture.  The glint of the sunlight on the Galilaean sea so near at
hand, with the uncounted flowers of the spring-time that covered
the lower plains, lent a charm to the scene that Quintus remembered
always.


At the outset the Roman convert is impressed with the goodly number
of those first disciples.  They are not twelve or six score, but
many more.  They greet each other with the salutation, "Peace be to
you," and then they rapturously add, "To-day we shall see our
Lord."  In that intimacy which should always mark the followers of
Christ, they give Quintus their welcome; and at once he feels
himself among a congenial brotherhood.

One is by name Nicodemus, a member of the Great Sanhedrin.  Another
is one Bartimaeus, from southern Jericho, whose finger tips have
been his eyes, till the Lord has healed his blindness.  A third has
been a demoniac among the hills of the Gergesenes, and has been a
wandering and truculent challenge to his times.  A woman is there
from Jacob's well, with Salome and Susanna and the virgin mother
herself.  They are from southern Bethlehem; they have come from the
wild hills of Peraea, beyond the Jordan; many are from Galilee,
where Christ has found so many devoted followers.  All these, as
well as the immortal eleven who have composed the inner circle of
the Master's associates.

Two other peculiar disciples does Quintus see, both of whom have
been raised from the dead.  Lazarus has come, who has so often
welcomed the Lord to his home in Bethany; and with him are the
sisters, of whom one has heard the Teacher say.  "Whosoever liveth
and believeth in me shall never die."  The other is a young
vineyard keeper from the neighboring village of Nain, whom Christ
has restored.  His word to Quintus is:

"Last year I sickened with a fever and passed through the door of
death.  They were carrying me out for burial, and my widowed mother
was weeping as one weeps who has lost her only son.  The Master
halted the mourners, and called me back to earth.  I have never
told of the wonders which I saw in the spirit world; it would not
be lawful.  But I have been in the great spaces beyond the stars,
and know that the tomb is only a resting place for a little sleep."

"How many disciples are there here?" Quintus asks of the good John.
To which question the other answers:

"Over a half thousand.  It has been our Master's wish that every
disciple of his throughout the land should come to this meeting
place.  Unto all he would show himself once more, before he returns
to the upper life.  So they shall have a glad memory of his face,
and shall be strengthened in their coming tribulations by the hope
of immortality."

Then suddenly--the risen Lord has come!  The marvel of it!  The
splendor of it!  While the five hundred are talking together, the
air grows luminous with his presence.  Out of the invisible he
appears.  As suddenly he comes as Aurora in her chariot drives up
the eastern sky and brings in the shining day.  When the company
have fallen on their faces and have adored their Master, in the
hush that follows he gives them a great commission:

"You are to go forth." he says, "and herald my gospel to the world.
Let there be no laggards in your company.  It is a lifelong charge.
There is a task for Petrus and Johannes, for Philippus and
Mattheus, and for all.  You are to look for disciples everywhere.
You are to proclaim the message of repentance.  You are to give
them the waters of baptism, in the name of the God triune.  You are
to declare to sad-hearted men the promise of eternal life, until I
shall come again to take men to myself."

That honorable commission!  It was in coming days to stir the souls
of apostles and quicken the feet of missioners and fire with zeal
earth's coming reformers.  Nor does Quintus forget that he too has
his charge.   In the city on the Tiber is to be his task.  To his
home circle, to priests in the temples of the gods, and even to the
royal Tiberius he is to herald the gospel of the resurrection.  His
vision of the risen Lord is the measure of his opportunity.

Then the Master looks into his very face, and remembers him as the
Roman knight he had seen in the Porch of Solomon.  The half
thousand disciples on Kurn Hattin prostrate themselves to the
earth; and in their acclaim the soldier joins his voice, "Rabboni!
Rabboni!  Our great Master!"  Then departs the Christ, and back to
their homes they go, evermore to comfort themselves with the vision
of their risen Lord.


Soon afterward their Rabboni goes from earth.  Out beyond the hill
of Olivet he walks one day with his eleven.  In their last words
together he reminds them again that they are to be his heralds to
the eastern world.  A cloud gathers above their heads, like some
halting chariot, and he is gone forever from human sight.  Yet only
in the distance it seems a cloud.  For John afterward says to
Quintus that it was in reality a phalanx of ten thousand angels,
robed in whiteness and sent to convoy the Son of God to glory
everlasting.

With Quintus that visit to Kurn Hattin shaped all his future.  His
Master's countenance had seemed to him more wonderful than any face
which the gifted Phidias had ever carved in stone.  But never in
after days could he worthily tell to Lucretia the vision he had
seen.  Only in one poor sentence could he sum it up: "I have seen
for myself the risen and ascending Lord."